\documentclass[a4paper,11pt]{article}

\usepackage[utf8]{inputenc}
\usepackage[T1]{fontenc}
\usepackage[ngerman]{babel}
\usepackage{amsmath}
\usepackage{amssymb}
\usepackage{xcolor}
\usepackage{colortbl}
\usepackage{array}
\usepackage{tikz}
\usepackage{needspace}
\usepackage{geometry}
\geometry{a4paper,margin=2.2cm}

% ----- Dark Mode (Pflicht) -----
\definecolor{bgdark}{HTML}{1E1E1E}   % dunkler Hintergrund
\definecolor{fgtext}{HTML}{E6E6E6}   % heller Text
\definecolor{accent}{HTML}{89B4FA}   % Akzent (hellblau)
\definecolor{accent2}{HTML}{F9E2AF}  % Akzent (gelb)
\definecolor{muted}{HTML}{A6ADC8}    % gedämpfter Text
\definecolor{gridcol}{HTML}{4A4A4A}  % gedämpfte Gitterfarbe
\definecolor{rulecol}{HTML}{6A6A6A}  % helle Linienfarbe
\pagecolor{bgdark}
\color{fgtext}

\setlength{\arrayrulewidth}{0.3pt}
\arrayrulecolor{rulecol}
\setlength{\parindent}{0pt}
\setlength{\parskip}{4pt}

% ----- Karo-Raster zum Ausfüllen -----
\newcommand{\karo}[1]{\par\noindent
  \begin{tikzpicture}
    \draw[step=5mm, gridcol, very thin] (0,0) grid (\linewidth,#1);
  \end{tikzpicture}\par}

% ----- Laufende Kopfzeile -----
\usepackage{fancyhdr}
\pagestyle{fancy}
\fancyhf{}
\renewcommand{\headrulewidth}{0pt}
\fancyhead[L]{\footnotesize\color{muted}\textsc{Numerische Methoden} $\cdot$ Pr\"ufung 01 Probeklausur}
\fancyhead[R]{\footnotesize\color{muted}Matr.-Nr.\ \textcolor{rulecol}{\rule{2.4cm}{0.3pt}}\quad \thepage}

% ----- Aufgabenkopf -----
\newcommand{\aufg}[3]{%
  \par\vspace{2pt}
  \noindent\textcolor{rulecol}{\rule{\linewidth}{0.3pt}}\\[3pt]
  {\large\bfseries\color{accent} Aufgabe #1\ \ \textcolor{fgtext}{\normalsize\mdseries #2}}\hfill
  {\bfseries\color{accent2}#3 Punkte}\\[1pt]
  \noindent\textcolor{rulecol}{\rule{\linewidth}{0.3pt}}
  \par\vspace{6pt}
}
\newcommand{\tp}[1]{\hfill{\small\color{muted}\textit{#1 Punkte}}}

\begin{document}

% =====================================================================
% DECKBLATT
% =====================================================================
\thispagestyle{fancy}

{\small\color{muted}\textsc{DHBW Ravensburg} \\ Duale Hochschule Baden-W\"urttemberg}

\vspace{10pt}
{\fontsize{34}{38}\selectfont\bfseries\color{accent} Pr\"ufung}\\[2pt]
{\itshape\color{fgtext} Numerische Methoden $\cdot$ Pr\"ufung 01 Probeklausur}

\vspace{10pt}
\noindent\textcolor{rulecol}{\rule{\linewidth}{0.4pt}}
\vspace{4pt}

\renewcommand{\arraystretch}{1.32}
\noindent\begin{tabular}{@{}p{4.2cm}p{11cm}@{}}
\textsc{\color{muted}Studiengang}     & W4DSKI\_108 Data Science und K\"unstliche Intelligenz\\
\textsc{\color{muted}Jahrgang}        & RV-WDS125 / RV-WDS225\\
\textsc{\color{muted}Modul}           & Fortgeschrittene Lineare Algebra und Analysis\\
\textsc{\color{muted}Veranstaltung}   & Numerische Methoden\\
\textsc{\color{muted}Pr\"ufer}        & Prof.\ Dr.-Ing.\ Mark Schutera\\
\textsc{\color{muted}Datum}           & \rule{4cm}{0.3pt}\\
\textsc{\color{muted}Bearbeitungszeit}& 90 Minuten\\
\textsc{\color{muted}Max.\ Punktzahl} & 75 Punkte\\
\textsc{\color{muted}Hilfsmittel}     & kein Taschenrechner. Open book; alle gedruckten oder
                                        handschriftlichen Unterlagen, die auf Papier mitgebracht werden.\\
\end{tabular}

\vspace{6pt}
\noindent\textcolor{rulecol}{\rule{\linewidth}{0.4pt}}

\vspace{10pt}
{\small\color{muted}\textsc{Anweisungen}}\par
\vspace{2pt}
{\small
Schreiben Sie Ihre L\"osungen und L\"osungswege leserlich in die vorgesehenen Bereiche.
Ordnen Sie jede L\"osung deutlich der zugeh\"origen Aufgabe zu. \textbf{Begr\"unden Sie Ihre
Antworten kurz: eine reine Zahl ohne Rechenweg gibt keine volle Punktzahl.} Br\"uche d\"urfen
als L\"osung stehen bleiben; Ableitungen d\"urfen analytisch gebildet werden, sofern nicht
ausdr\"ucklich anders verlangt. Lesen Sie die Aufgabenstellung sorgf\"altig, bevor Sie mit der
Bearbeitung beginnen. Viel Erfolg!}

\vspace{12pt}
{\small\color{muted}\textsc{Durch Pr\"ufer auszuf\"ullen}}\par
\vspace{3pt}
\renewcommand{\arraystretch}{1.4}
\noindent\begin{tabular}{|l|c|c|c|c|c|c|c|}
\hline
\color{muted}Aufgabe & 1 & 2 & 3 & 4 & 5 & 6 & $\Sigma$\\
\hline
\color{muted}Maximal & 16 & 7 & 14 & 13 & 14 & 11 & 75\\
\hline
\color{muted}Erreicht & & & & & & & \\
\hline
\end{tabular}

% =====================================================================
% AUFGABE 1 -- Block 1
% =====================================================================
\clearpage
\aufg{1}{Gleitkomma, Festkomma \& Brute Force}{16}

\textbf{(a)} Sie entwerfen ein Gleitkommasystem (floating-point system) mit $1$ Vorzeichenbit
und weiteren Bits, die Sie frei auf Mantisse (mantissa) und Exponent (exponent) verteilen.
\tp{6}

\smallskip
Diskutieren Sie den Trade-off: Welche Eigenschaft verbessert sich, wenn Sie ein Bit der
\emph{Mantisse} zuweisen, welche, wenn Sie es dem \emph{Exponenten} geben? Argumentieren Sie
\"uber Aufl\"osung (resolution) und Wertebereich (dynamic range).

\smallskip
Fixieren wir die Aufteilung auf $3$ Mantissenbits (normalisierte Mantisse $1{.}b_1b_2b_3$ mit
implizitem f\"uhrenden Bit) und $3$ Exponentenbits mit Bias $3$. Bestimmen Sie das
Maschinenepsilon $\varepsilon_{\mathrm{mach}}$ und berechnen Sie in diesem System
\[
   \frac{1}{\bigl(1+\tfrac{1}{32}\bigr)-1}.
\]
Geben Sie den Nenner nach Rundung sowie das Endergebnis an, vergleichen Sie mit dem
mathematischen Wert und erkl\"aren Sie die Ausl\"oschung (catastrophic cancellation).

\karo{7.5cm}

\clearpage
\textbf{(b)} Ein \emph{Festkommasystem} (fixed-point) verwende $8$ Bit, davon $4$ Vorkomma-
und $4$ Nachkommabits (Skalierungsfaktor $2^{-4}$).
\tp{3}

\smallskip
Geben Sie die Aufl\"osung und die gr\"o\ss te darstellbare (vorzeichenlose) Zahl an.
Stellen Sie $3{,}25$ in diesem System bin\"ar dar. Vergleichen Sie Festkomma und Gleitkomma:
Wie unterscheiden sich die Abst\"ande benachbarter darstellbarer Zahlen, und welche Rolle
spielt der Skalierungsfaktor?

\karo{7cm}

\clearpage
\textbf{(c)} Wir diskretisieren das Intervall $[0,1]$ mit $N=4$ gleich gro\ss en
Teilintervallen.
\tp{3}

\smallskip
Berechnen Sie die Schrittweite $h=\frac{b-a}{N}$ und geben Sie alle Gitterpunkte an.
Ordnen Sie anschlie\ss end die folgenden drei Fehlerquellen jeweils einem Fehlertyp zu
(\emph{Modellfehler}, \emph{Abschneidefehler}, \emph{Rundungsfehler}) und begr\"unden Sie kurz:
\begin{enumerate}\setlength{\itemsep}{1pt}
\item[(i)]   Bei der Simulation wird die Luftreibung vernachl\"assigt.
\item[(ii)]  Die Reihe $e^x=\sum_{n\ge0}\tfrac{x^n}{n!}$ wird nach drei Gliedern abgebrochen.
\item[(iii)] Der Wert $\tfrac13$ wird als $0{,}333$ gespeichert.
\end{enumerate}

\karo{6.5cm}

\clearpage
\textbf{(d)} Wir suchen die L\"osung des linearen Systems
\[
  \begin{pmatrix} 2 & 1 \\ 1 & 3 \end{pmatrix}
  \begin{pmatrix} w \\ b \end{pmatrix}
  =
  \begin{pmatrix} 3 \\ 5 \end{pmatrix}
\]
mit einer Gitter-Suche (grid search) \"uber $w,b\in\{0,1,2\}$.
\tp{4}

\smallskip
W\"ahlen Sie ein Fehlerma\ss\ und begr\"unden Sie es kurz; bestimmen Sie $w^\ast,b^\ast$.
Vergleichen Sie mit der exakten L\"osung. Findet die Gitter-Suche das System exakt? Wie l\"asst
sich die Methode verbessern?

\smallskip
{\small\color{muted}\textsc{Hinweis.} Exakte L\"osung: $w^\star=4/5$, $b^\star=7/5$.}

\karo{8cm}

% =====================================================================
% AUFGABE 2 -- Block 2
% =====================================================================
\clearpage
\aufg{2}{Fehleranalyse}{7}

Wir bestimmen die Ableitung $f'$ von $f(x)=x^3$ (wahres Problem: $f'(x)=3x^2$) numerisch mit
der Vorw\"artsdifferenz
\[
  D_h(x)=\frac{f(x+h)-f(x)}{h}.
\]
Die Eingabe sei durch einen Messfehler um $\tilde x=x+\delta$ gest\"ort. Untersuchen Sie
nacheinander die vier Eigenschaften der Fehleranalyse.
\tp{7}

\smallskip
\emph{Kondition (conditioning):} Quantifizieren Sie die Kondition von $f'(x)=3x^2$ an den
Stellen $x=0$ und $x=2$; benennen Sie den gut bzw.\ schlecht konditionierten Bereich.

\smallskip
\emph{Konsistenz (consistency):} Zeigen Sie $D_h(x)=3x^2+3xh+h^2$ und quantifizieren Sie den
Konsistenzfehler $|D_h(x)-f'(x)|$ am Punkt $x=1$ f\"ur $h_1=1$ und $h_2=\tfrac12$. Welche
Schrittweite ist konsistenter?

\smallskip
\emph{Stabilit\"at (stability) \& Konvergenz (convergence):} Erkl\"aren Sie, warum eine
\emph{immer kleinere} Schrittweite bei endlicher Rechengenauigkeit nicht zwingend zu einer
besseren Approximation f\"uhrt (Konsistenz- vs.\ Rundungsfehler).

\karo{9cm}

% =====================================================================
% AUFGABE 3 -- Block 3 + 4
% =====================================================================
\clearpage
\aufg{3}{Nullstellen \& Globale Optimierung}{14}

\textbf{(a)} Wir suchen eine Nullstelle (root) von $f(x)=x^2-3$.
\tp{8}

\begin{center}
\begin{tikzpicture}[scale=0.8]
  \draw[->,rulecol] (-1.4,0) -- (3.2,0) node[right]{\small$x$};
  \draw[->,rulecol] (0,-3.6) -- (0,6.4) node[above]{\small$f(x)$};
  \foreach \y in {3,6} \draw[rulecol] (-0.08,\y) -- (0.08,\y) node[left=2pt]{\scriptsize$\y$};
  \foreach \x in {1,2,3} \draw[rulecol] (\x,-0.08) -- (\x,0.08) node[below=2pt]{\scriptsize$\x$};
  \draw[accent,thick,domain=-1.3:3.05,samples=80,smooth] plot (\x,{(\x)*(\x)-3});
  \draw[accent2,thick] (2,1) circle (2.5pt);
  \node[accent2] at (2.45,1.4) {\scriptsize$x_0$};
\end{tikzpicture}
\end{center}

\textbf{Leiten Sie} das Newton-Verfahren $x_{n+1}=x_n-\frac{f(x_n)}{f'(x_n)}$ \emph{aus der
Taylor-Entwicklung} erster Ordnung her. F\"uhren Sie dann mit $x_0=2$ zwei Iterationen durch
($x_1,x_2$ als Bruch) und nennen Sie das Konvergenzkriterium.

\smallskip
Nehmen Sie an, $f$ sei nur durch Funktionsauswertungen zug\"anglich. Geben Sie eine numerische
Approximation f\"ur $f'(x_0)$ samt Approximationsfehler an und leiten Sie daraus ein
\emph{ableitungsfreies} Verfahren ab, das einen Versagensmodus von Newton umgeht. Benennen Sie
diesen Versagensmodus.

\karo{8.5cm}

\clearpage
\textbf{(b)} Wir wechseln zur Minimumsuche.
\tp{6}

\smallskip
\emph{(i)} Erkl\"aren Sie kurz den Unterschied zwischen \emph{lokalem} und \emph{globalem}
Minimum und warum eine multimodale Landschaft wie die \emph{Shekel-Foxholes} (viele
``L\"ocher'') f\"ur lokale Optimierer schwierig ist.

\smallskip
\emph{(ii)} F\"ur die Funktion $g(x)=x^3-3x$ f\"uhren Sie \emph{einen} Schritt des
Newton-Verfahrens \emph{f\"ur Optima}
\[
   x_{n+1}=x_n-\frac{g'(x_n)}{|g''(x_n)|}
\]
mit Startwert $x_0=2$ durch.

\smallskip
\emph{(iii)} Bei \emph{zuf\"alligen Neustarts} (random restarts) treffe ein Start mit
Wahrscheinlichkeit $p=\tfrac12$ das globale Becken. Berechnen Sie $P=1-(1-p)^K$ f\"ur $K=3$
und bestimmen Sie das kleinste $K$ mit $P>0{,}9$.

\smallskip
\emph{(iv)} Nennen Sie einen weiteren Ansatz gegen raue Landschaften (z.\,B.\ Gl\"attung,
Basin Hopping, Simulated Annealing) und w\"ahlen Sie ihn passend.

\karo{8.5cm}

% =====================================================================
% AUFGABE 4 -- Block 5
% =====================================================================
\clearpage
\aufg{4}{Differenzialgleichungen}{13}

Anfangswertproblem $y'=t+y$, $\;y(0)=1$.

\smallskip
\textbf{(a)} Berechnen Sie zwei Schritte des expliziten Euler-Verfahrens mit $h=1$
(N\"aherungen f\"ur $y(1)$ und $y(2)$).
\tp{3}

\karo{5cm}

\clearpage
\textbf{(b)} Gegeben das Butcher-Tableau
\(
  \renewcommand{\arraystretch}{1.2}
  \begin{array}{c|cc} 0 & 0 & 0\\ \tfrac12 & \tfrac12 & 0\\ \hline & 0 & 1\end{array}
\).
Bestimmen Sie die Ordnung (\"uber $\sum_i b_i=1$ und $\sum_i b_i c_i=\tfrac12$) und f\"uhren Sie
\emph{einen} Schritt dieses Runge-Kutta-Verfahrens ($k_1,k_2,y_1$) f\"ur das Anfangswertproblem
mit $h=1$ aus.
\tp{4}

\karo{7.5cm}

\clearpage
\textbf{(c)} F\"uhren Sie \emph{einen} Schritt des klassischen \textbf{RK4}-Verfahrens
\[
  y_{n+1}=y_n+\tfrac{h}{6}\bigl(k_1+2k_2+2k_3+k_4\bigr)
\]
mit $h=1$ f\"ur dasselbe Anfangswertproblem aus. Berechnen Sie $k_1,k_2,k_3,k_4$ und $y(1)$ und
vergleichen Sie mit dem Euler- und dem RK2-Ergebnis.
\tp{4}

\karo{8cm}

\clearpage
\textbf{(d)} Betrachten Sie $y'=-\lambda y$, $\lambda>0$. F\"ur welche Schrittweiten $h$ ist das
explizite Euler-Verfahren stabil? Geben Sie die Stabilit\"atsbedingung an und beschreiben Sie,
was f\"ur zu gro\ss e $h$ passiert.
\tp{2}

\karo{5.5cm}

% =====================================================================
% AUFGABE 5 -- Block 6
% =====================================================================
\clearpage
\aufg{5}{Integration \& Interpolation}{14}

\textbf{(a)} Bestimmen Sie das Lagrange-Interpolationspolynom $p(x)$ durch
$(0,1),(1,3),(2,7)$ und werten Sie es bei $x=\tfrac12$ aus. Erkl\"aren Sie au\ss erdem kurz,
was ein \emph{(kubischer) Spline} ist und welchen Vorteil er gegen\"uber einem
Interpolationspolynom hohen Grades hat (Stichwort Runge-Ph\"anomen).
\tp{4}

\karo{7.5cm}

\clearpage
\textbf{(b)} Wir betrachten $I=\int_0^2 (x^2+1)\,dx$ mit St\"utzwerten $f(0)=1,f(1)=2,f(2)=5$.
\tp{5}

\smallskip
Geben Sie die Gewichte/Formeln der Newton-Cotes-Regeln \emph{Mittelpunkt}, \emph{Trapez},
\emph{Simpson}, \emph{3/8} und \emph{Boole} an. Berechnen Sie $I$ dann mit der Trapezregel
($h=1$), der Simpson-Regel ($h=1$) und der 3/8-Regel (3 Teilintervalle, $h=\tfrac23$),
vergleichen Sie mit dem exakten Wert und erkl\"aren Sie, warum Simpson und 3/8 hier exakt sind,
die Trapezregel aber nicht.

\karo{8.5cm}

\clearpage
\textbf{(c)} Die Romberg-Methode extrapoliert die Trapezregel (Richardson). Mit den
Trapezwerten $T(h)$ und $T(h/2)$ lautet der erste Romberg-Schritt
\[
  R = \frac{4\,T(h/2)-T(h)}{3}.
\]
Berechnen Sie $T(h{=}2)$ und $T(h{=}1)$ f\"ur $I$ aus (b) und daraus $R$. Was f\"allt auf?
\tp{2}

\karo{6cm}

\clearpage
\textbf{(d)} Erkl\"aren Sie die Monte-Carlo-Integration mit Sch\"atzer
$\hat I=|\Omega|\,\tfrac1N\sum_{i=1}^N f(X_i)$ und Standardfehler
$\mathrm{SE}=|\Omega|\,\sigma_f/\sqrt N$. Um welchen Faktor muss $N$ wachsen, um den SE zu
halbieren? Warum ist Monte-Carlo im Hochdimensionalen vorteilhaft (Fluch der Dimensionen)?
\tp{3}

\karo{6.5cm}

% =====================================================================
% AUFGABE 6 -- Block 7
% =====================================================================
\clearpage
\aufg{6}{Fourier-Transformation}{11}

\textbf{(a)} Berechnen Sie f\"ur das diskrete Signal $x=[\,1,0,-1,0\,]$ ($N=4$) die DFT
\[
   X_k=\sum_{n=0}^{3} x_n\,e^{-i\,2\pi k n/4},\qquad k=0,1,2,3,
\]
von Hand ($e^{-i\pi/2}=-i$). Geben Sie $|X_k|$ an und interpretieren Sie das Spektrum.
\tp{5}

\karo{8cm}

\clearpage
\textbf{(b)} Geben Sie die \emph{kontinuierliche} Fourier-Transformation $F(\omega)$ und ihre
Rücktransformation an. Betrachten Sie das Interferenzsignal
\[
   f(t)=A_1\cos(\omega_1 t)+A_2\cos(\omega_2 t),\qquad \omega_1\neq\omega_2 .
\]
Wie viele Peaks zeigt $|F(\omega)|$ und bei welchen Frequenzen? Nutzen Sie die Euler-Formel
$e^{i\omega t}=\cos(\omega t)+i\sin(\omega t)$ zur Begr\"undung.
\tp{4}

\karo{7.5cm}

\clearpage
\textbf{(c)} Ein Sensor tastet mit $f_s=8\,\mathrm{Hz}$ ab und liefert $N=8$ Werte. Im
Betragsspektrum liegt ein Peak im Bin $k=2$. Welcher Frequenz $f$ entspricht das
($f_k=k\,f_s/N$)? Bis zu welcher Frequenz ist eindeutige Darstellung m\"oglich (Nyquist), und
was passiert mit h\"oheren Frequenzen?
\tp{2}

\karo{5.5cm}

\end{document}
